%------------------------------------------------------------------------------%
\documentclass[fleqn,utf8,aspectratio=169,12pt]{beamer}

\usepackage{integracao_e_series}

\title{Integrais de Funções Trigonométricas}
\subtitle{3 -- Cotangente e Cossecante}

%------------------------------------------------------------------------------%
\begin{document}

\addtitle

%------------------------------------------------------------------------------%
\begin{frame}{Cotangente e Cossecante}
  Mesmas estratégias
  \[
    \int \cot^m(x) \csc^m(x)\,dx
  \]
\end{frame}

%% 1
%%------------------------------------------------------------------------------%
%\addtocounter{exemplo}{1}
%\section{Exemplo \theexemplo}
%
%%------------------------------------------------------------------------------%
%\begin{frame}{Exemplo \theexemplo}
%\end{frame}
%
%
%%------------------------------------------------------------------------------%
%\section{Lista Mínima}
%
%%------------------------------------------------------------------------------%
%\begin{frame}{Lista Mínima}
%  Estudar a Seção 4.3 da Apostila
%
%  Exercícios: 1i, 1j, 1m
%
%  \vfill
%  Atenção: A prova é baseada no livro, não nas apresentações
%\end{frame}

%------------------------------------------------------------------------------%
\end{document}
%------------------------------------------------------------------------------%
